% ============================================================
% Section 95: 德鲁德模型电导率与铜的弛豫时间
% ============================================================
\section{固体理论：德鲁德模型电导率与铜的弛豫时间}
\label{sec:drude-conductivity}

\begin{modelbox}[题目：金属电导率公式的推导与铜中电子弛豫时间]
在金属自由电子模型中，推导金属电导率
\[
\sigma=\frac{ne^2\tau}{m},
\]
并估计铜中电子的弛豫时间 $\tau$。铜的电阻率为 $1.7\times10^{-6}\,\Omega\cdot\text{cm}$，原子密度为 $8.5\times10^{22}\,\text{cm}^{-3}$。
\end{modelbox}

\begin{tipbox}[考点快速分析]
\begin{itemize}
    \item \textbf{半经典运动}：$\hbar\,dk/dt=F=-e\mathcal{E}$
    \item \textbf{平均自由时间}：两次碰撞间电子获得漂移速度 $v_d=-e\mathcal{E}\tau/m$
    \item \textbf{欧姆定律}：$j=-nev_d=\sigma\mathcal{E}$
    \item \textbf{弛豫时间}：$\tau=m/(ne^2\rho)$，铜中约 $10^{-14}\,$s
\end{itemize}
\end{tipbox}

\begin{derivationbox}[标准解题过程]
\textbf{Step 1: 电导率公式的推导}

电子在外电场 $\mathcal{E}$ 中受力 $F=-e\mathcal{E}$。半经典运动方程
\begin{equation}
\hbar\frac{d\bm{k}}{dt}=-e\bm{\mathcal{E}}.
\end{equation}

在两次碰撞之间的平均时间 $\tau$ 内，电子波矢的增量
\begin{equation}
\delta\bm{k}=\frac{F}{\hbar}\tau=-\frac{e\bm{\mathcal{E}}}{\hbar}\tau.
\end{equation}

相应的漂移速度
\begin{equation}
\bm{v}_d=\frac{\hbar\delta\bm{k}}{m}=-\frac{e\bm{\mathcal{E}}}{m}\tau.
\end{equation}

电流密度（$n$ 个电子，每个电荷 $-e$）
\begin{equation}
\bm{j}=-ne\bm{v}_d=-ne\Bigl(-\frac{e\tau}{m}\bm{\mathcal{E}}\Bigr)=\frac{ne^2\tau}{m}\bm{\mathcal{E}}.
\end{equation}

由欧姆定律 $\bm{j}=\sigma\bm{\mathcal{E}}$，得\textbf{德鲁德电导率公式}
\begin{equation}
\boxed{\sigma=\frac{ne^2\tau}{m}.}
\end{equation}

\textbf{Step 2: 铜中电子弛豫时间的估算}

铜的电阻率
\begin{equation}
\rho=1.7\times10^{-6}\,\Omega\cdot\text{cm}=1.7\times10^{-8}\,\Omega\cdot\text{m}.
\end{equation}

每个铜原子贡献 1 个自由电子，电子数密度
\begin{equation}
n=8.5\times10^{22}\,\text{cm}^{-3}=8.5\times10^{28}\,\text{m}^{-3}.
\end{equation}

由 $\sigma=1/\rho=ne^2\tau/m$，反解
\begin{equation}
\tau=\frac{m}{ne^2\rho}.
\end{equation}

代入数值：
\begin{align}
\tau&=\frac{9.11\times10^{-31}}{8.5\times10^{28}\times(1.60\times10^{-19})^2\times1.7\times10^{-8}} \notag\\
&=\frac{9.11\times10^{-31}}{8.5\times2.56\times1.7\times10^{28-38-8}} \notag\\
&=\frac{9.11\times10^{-31}}{36.99\times10^{-18}} \notag\\
&\approx\boxed{2.5\times10^{-14}\,\text{s}.}
\end{align}
\end{derivationbox}

\begin{formulabox}[答案汇总]
\begin{center}
\begin{tabular}{ll}
\toprule
\textbf{项目} & \textbf{表达式 / 数值} \\
\midrule
电导率公式 & $\sigma=ne^2\tau/m$ \\
铜的弛豫时间 & $\tau\approx2.5\times10^{-14}\,\text{s}$ \\
\bottomrule
\end{tabular}
\end{center}
\end{formulabox}

\begin{insightbox}[物理图像]
\begin{itemize}
    \item 德鲁德模型通过引入平均弛豫时间 $\tau$，成功将微观电子散射过程与宏观电导率联系起来。$\tau$ 是电子在两次散射之间自由飞行时间的平均值。
    \item 铜的 $\tau\sim10^{-14}\,$s，若结合费米速度 $v_F\sim1.6\times10^6\,$m/s，可得室温平均自由程 $l=v_F\tau\sim400\,$\AA。这一长度远大于晶格常数（$\sim2\,$\AA），说明电子在晶格中传播时``感受''到的是平均势场而非单个原子势——这正是能带论和准经典近似的物理基础。
    \item 电阻率 $\rho\propto1/\tau$，而 $\tau$ 主要由电子-声子散射和电子-杂质散射决定。温度升高时声子增多，$\tau$ 减小，电阻率增大——这解释了金属电阻的正温度系数。
\end{itemize}
\end{insightbox}
